\chapter{Conclusiones}
\label{cap:conclusiones}

Este capítulo recapitula el trabajo realizado, valora el grado de consecución de los
objetivos planteados, expone las aportaciones y limitaciones, y traza las líneas de trabajo
futuro. Por último, se incluye una breve valoración personal.

%% =============================================================================
\section{Grado de consecución de los objetivos}
\label{sec:consecucion-objetivos}

El desarrollo de este Trabajo Fin de Grado ha permitido construir un sistema robótico de
patrulla autónoma funcional e integrado. Retomando los objetivos específicos formulados en
el Capítulo~\ref{cap:objetivos}, puede afirmarse que se ha implementado la cartografía
autónoma del entorno mediante exploración y \gls{slam} (OE1) y la generación y ejecución de
una ruta de cobertura con \gls{nav2} (OE2). El ciclo operativo se coordina mediante un
orquestador en forma de máquina de estados que encadena exploración, mapeado, planificación,
patrulla y retorno al punto de carga (OE3). Sobre esa base, se ha integrado la detección de
personas como incidencia con \gls{yolo} (OE4), se ha diseñado la arquitectura orientada a
eventos que desacopla el robot del backend (OE5) y se ha desarrollado el servicio de backend
con persistencia de datos relacionales, de vídeo y de series temporales (OE6). Asimismo, se
ha implementado el esquema de seguridad basado en \gls{tofu}, \gls{jwt} y control de acceso
por roles (OE7) y se ha construido la aplicación de usuario multiplataforma con
\mbox{Flutter} (OE8). Por último, la integración de todos los subsistemas en una solución
desplegable mediante contenedores se ha completado en su vertiente funcional, validada de
extremo a extremo en simulación (OE9), si bien la evaluación cuantitativa sobre hardware
real queda pendiente.

%% =============================================================================
\section{Aportaciones}
\label{sec:aportaciones}

La principal aportación del trabajo es la \textbf{integración de extremo a extremo} de un
sistema de vigilancia robótica con tecnologías de código abierto, que cubre desde la
cartografía y la patrulla autónoma hasta la notificación de incidencias y su consulta desde
una aplicación multiplataforma. Frente a soluciones centradas en un único subproblema, este
trabajo articula navegación, percepción, mensajería, servicios web y aplicación de usuario
en un único sistema coherente y desplegable. Más allá de esa integración, pueden destacarse
varias aportaciones concretas:

\begin{itemize}
  \item \textbf{Arquitectura orientada a eventos.} El desacoplamiento de la capa robótica y
    el backend mediante un bus de mensajería ha demostrado su valor al permitir desarrollar,
    probar y desplegar cada componente de forma independiente, y al absorber la
    indisponibilidad temporal de cualquiera de ellos sin pérdida de eventos.
  \item \textbf{Orquestación autónoma del ciclo de patrulla.} El diseño del orquestador como
    una máquina de estados que encadena exploración, mapeado, planificación y patrulla
    encapsula en un único nodo la coordinación de capacidades heterogéneas, y reutiliza el
    mapa cuando ya existe para no reexplorar el entorno innecesariamente.
  \item \textbf{Modelo de seguridad para el robot.} La autenticación del robot mediante el
    patrón \gls{tofu} y tokens \gls{jwt}, con doble verificación de identidad (en las
    dependencias del backend y, de nuevo, en el consumidor de eventos), aporta un esquema de
    confianza adecuado a un dispositivo que opera sin supervisión continua.
  \item \textbf{Despliegue reproducible.} La descripción declarativa de la infraestructura
    mediante contenedores permite levantar el sistema completo de forma uniforme, lo que
    facilita tanto el desarrollo como una futura puesta en producción.
\end{itemize}

%% =============================================================================
\section{Limitaciones}
\label{sec:limitaciones}

El trabajo presenta varias limitaciones que conviene reconocer. La más relevante es que la
validación realizada es funcional y se ha llevado a cabo en simulación, sin una campaña de
medición cuantitativa sobre hardware real que respalde el rendimiento del sistema; las
métricas de navegación, detección y latencia quedan, por tanto, sin caracterizar. Además, la
estrategia de planificación de la ruta de cobertura es deliberadamente sencilla (un barrido
en zig-zag) y puede resultar subóptima en estancias de geometría compleja. La detección se
limita actualmente a la presencia de personas, sin seguimiento de objetos ni distinción
entre individuos, lo que puede generar alertas redundantes ante una misma persona.

En el plano de la seguridad deben reconocerse dos limitaciones conocidas: el patrón
\gls{tofu} asume que el primer contacto del robot con el backend es legítimo, asunción que
sería explotable si un atacante se adelantara al aprovisionamiento; y la política de
\gls{cors} permanece abierta. Ambas resultan asumibles en el despliegue local y controlado
actual, pero deben endurecerse antes de un uso en producción. Por último, algunas
capacidades que se contemplaron en un principio no forman parte de la implementación actual,
en concreto la transmisión de vídeo en directo, el seguimiento multiobjeto y las
notificaciones en tiempo real; todas ellas se recogen como líneas de trabajo futuro.

%% =============================================================================
\section{Trabajo futuro}
\label{sec:trabajo-futuro}

A partir de las limitaciones anteriores se identifican las siguientes líneas de trabajo
futuro:

\begin{itemize}
  \item \textbf{Evaluación cuantitativa sobre hardware real.} Ejecutar una campaña de
    pruebas sistemática sobre el robot real que permita obtener las métricas previstas en el
    Capítulo~\ref{cap:resultados} y recalibrar los parámetros ajustados en simulación.
  \item \textbf{Endurecimiento de la seguridad.} Restringir la política de \gls{cors} y
    reforzar el aprovisionamiento inicial del robot frente al riesgo del primer contacto, de
    cara a un despliegue en producción.
  \item \textbf{Seguimiento multiobjeto.} Incorporar un algoritmo de seguimiento que asigne
    identificadores persistentes a las personas detectadas y reduzca las alertas
    redundantes.
  \item \textbf{Transmisión de vídeo en directo.} Añadir la difusión del vídeo en tiempo
    real para permitir la supervisión remota durante la patrulla.
  \item \textbf{Notificaciones en tiempo real.} Avisar al usuario de forma inmediata ante
    una incidencia mediante notificaciones \textit{push} o canales de comunicación en tiempo
    real.
  \item \textbf{Planificación de cobertura avanzada.} Sustituir la rejilla en zig-zag por
    estrategias de cobertura más eficientes y adaptadas a la geometría del entorno.
  \item \textbf{Operación multirobot.} Extender el sistema para coordinar varios robots
    patrullando un mismo espacio.
\end{itemize}

%% =============================================================================
\section{Valoración personal}
\label{sec:valoracion-personal}

La realización de este trabajo me ha permitido integrar en un único proyecto conocimientos
de áreas muy diversas (robótica, visión por computador, sistemas distribuidos y desarrollo
de aplicaciones) y enfrentarme a los retos que surgen al hacer que componentes heterogéneos
colaboren de forma coherente. Construir el sistema de extremo a extremo, y no solo una de
sus partes, ha sido la faceta más exigente y, a la vez, la más enriquecedora del trabajo.

Más allá del resultado técnico, considero que la principal lección ha sido la importancia de
un buen diseño arquitectónico: las decisiones tomadas al principio, en especial el
desacoplamiento mediante eventos, condicionaron favorablemente todo el desarrollo posterior
y facilitaron cada incorporación sucesiva. Llevar el sistema hasta su validación sobre
hardware real y completar su evaluación cuantitativa es el paso natural con el que me
gustaría dar continuidad a lo aquí construido.
